\documentclass[12pt]{article}
\usepackage{float}
\usepackage{graphicx}
\usepackage{fontspec}
\usepackage{url}
\usepackage{apacite}
\setmainfont{Times New Roman}
\usepackage[letterpaper,top=2cm,bottom=2cm,left=2cm,right=2cm,marginparwidth=1.5cm]{geometry}

\title{Weekly Assignment 5: Reflection on information security policy crafting}
\author{Truong Khang Tobias Nguyen}
\date{}

\begin{document}

\maketitle

\section{Learnings}
There are a plethora of things that I have learned from the NordBank case that we have discussed in class.
\newline
The most important thing is that developing and implementing appropriate security policies is not as straightforward as it sounds like. The NordBank case demonstrated that policy development inside of an organization involves more variables than I have initially expected such as the roles of management and departments, company hierarchies and security culture. In that sense, organizational factors can impact the development of security policies even though they might not be related on a surface level \cite{balagopal2024exploring}.
\newline
Subsequently, developing security policies involves several layers and factors that need to be considered in order to be properly developed and implemented. This was seen in the amount of people I would have invited to participate in the process, such as the CEO, CIO, the VPs and the legal and compliance team. 

\section{Reflections}
The NordBank case has challenged my previous assumptions about security policy development in ways that I have a more grounded view on it as my initial view on security policy development so far has been relatively straightforward up to this point. 
\newline
I initially believed that an organizations security departments would formulate the policies and that the rest of organization would just have to take it. During the course of the lecture, I came to the conclusion that such a procedure might lead to weak acceptance and non-compliance \cite{d2019predicting} and that strategies need to be employed to establish acceptance and legitimacy.
\newline
Additionally, I also believed that security departments developed and iterated on security policies as time passes and technology advances and have the rest of the organization follow the policies as non-compliance will lead to consequences. In that sense, the rest of the organization took the policies at face value and would not dare to break the rules out of fear of repercussions which maps pretty well to the concept of deterrence theory \cite{jervis1979deterrence}.
\newline
That being said, I have realized that developing security policies inside of an organization requires not only meaningful measures, but also needs to seek approval from multiple components inside of such a system. As an organization does not only involve the security department, and the security department not directly contributing to daily operations or revenue, they can not just simply craft policies and expect the rest to follow them, but this process requires the people closest to operations and management to participate and approve of them in order to develop meaningful policies that are accepted and complied with \cite{flowerday2016information}.

\bibliographystyle{apacite}
\bibliography{references}

\end{document}